\documentclass{article}
\usepackage{amsmath,amssymb,amsthm}
\usepackage{algorithm}
\usepackage{algpseudocode}
\usepackage{hyperref}
\usepackage{bm}
\usepackage{enumitem}
\newtheorem{theorem}{Theorem}[section]
\newtheorem{lemma}{Lemma}[section]
\newtheorem{assumption}{Assumption}[section]
\newtheorem{corollary}{Corollary}[section]

\begin{document}

%%%%%%%%%%%%%%%%%%%%%%%%%%%%%%%%%%%%%%%%%%%%%%%%%%%%%%%%%%%%%%%%%%%%%%%%%%%%%
\section*{Primal–Dual Hybrid Gradient (PDHG)}
%
The PDHG method solves a convex–concave saddle‑point problem
\[
\min_{x\in\mathbb{R}^n}\; \max_{y\in\mathbb{R}^m}
\; \mathcal{L}(x,y)\;:=\; f(x)+\langle Kx,y\rangle - g^{*}(y),
\tag{P}
\]
where
\begin{itemize}
  \item $f:\mathbb{R}^n\to\mathbb{R}$ is convex, differentiable and $L$‑smooth,
  \item $g:\mathbb{R}^m\to\mathbb{R}\cup\{+\infty\}$ is proper, convex and lower‑semicontinuous,
  \item $g^{*}$ denotes the convex conjugate of $g$,
  \item $K\in\mathbb{R}^{m\times n}$ is a linear operator.
\end{itemize}
The algorithm maintains a primal iterate $x_k$, a dual iterate $y_k$ and an extrapolated primal point
$\bar{x}_k$.

%%%%%%%%%%%%%%%%%%%%%%%%%%%%%%%%%%%%%%%%%%%%%%%%%%%%%%%%%%%%%%%%%%%%%%%%%%%%%
\section*{Mathematical Formulation}
Given step sizes $\tau>0$, $\sigma>0$ and an extrapolation parameter $\theta\in[0,1]$,
the PDHG updates are
\begin{align}
    y_{k+1} &= \operatorname{prox}_{\sigma g^{*}}\!\bigl(y_k+\sigma K\bar{x}_k\bigr),
    \label{eq:dual-update}\\[4pt]
    x_{k+1} &= \operatorname{prox}_{\tau f}\!\bigl(x_k-\tau K^{\!\top} y_{k+1}\bigr),
    \label{eq:primal-update}\\[4pt]
    \bar{x}_{k+1} &= x_{k+1}+\theta\bigl(x_{k+1}-x_k\bigr).
    \label{eq:extrapolation}
\end{align}
The proximal operator of a proper convex function $h$ is
\[
\operatorname{prox}_{\lambda h}(z)
:=\arg\min_{u}\Bigl\{h(u)+\frac{1}{2\lambda}\|u-z\|^{2}\Bigr\}.
\]

%%%%%%%%%%%%%%%%%%%%%%%%%%%%%%%%%%%%%%%%%%%%%%%%%%%%%%%%%%%%%%%%%%%%%%%%%%%%%
\section*{Pseudocode}
\begin{algorithm}[H]
\caption{Primal–Dual Hybrid Gradient}
\begin{algorithmic}[1]
\Require $f$, $g$, linear operator $K$, step sizes $\tau,\sigma>0$, extrapolation $\theta\in[0,1]$, 
initial $x_0\in\mathbb{R}^n$, $y_0\in\mathbb{R}^m$
\State $\bar{x}_0 \gets x_0$
\For{$k=0,1,2,\dots$}
    \State $y_{k+1}\gets\operatorname{prox}_{\sigma g^{*}}\bigl(y_k+\sigma K\bar{x}_k\bigr)$ \Comment{dual step}
    \State $x_{k+1}\gets\operatorname{prox}_{\tau f}\bigl(x_k-\tau K^{\!\top}y_{k+1}\bigr)$ \Comment{primal step}
    \State $\bar{x}_{k+1}\gets x_{k+1}+\theta\bigl(x_{k+1}-x_k\bigr)$ \Comment{extrapolation}
\EndFor
\State \textbf{output} $(x_{K},y_{K})$ for a chosen stopping index $K$
\end{algorithmic}
\end{algorithm}

%%%%%%%%%%%%%%%%%%%%%%%%%%%%%%%%%%%%%%%%%%%%%%%%%%%%%%%%%%%%%%%%%%%%%%%%%%%%%
\section*{Dry Run Example}
Consider the scalar problem
\[
\min_{x\in\mathbb{R}} \; f(x)= (x-3)^{2},
\]
which fits \eqref{eq:primal-update} with $K=1$, $g\equiv0$ (hence $g^{*}=0$ and $\operatorname{prox}_{\sigma g^{*}}$ is the identity).  
Take $\tau=0.1$, $\sigma=0.1$, $\theta=0$, $x_{0}=0$, $y_{0}=0$.

\begin{center}
\begin{tabular}{c|c|c|c}
$k$ & $y_{k+1}$ (dual) & $x_{k+1}$ (primal) & $f(x_{k+1})$ \\ \hline
0 &
$y_{1}=y_{0}+\sigma\bar{x}_{0}=0+0.1\cdot0=0$ &
$x_{1}= \operatorname{prox}_{\tau f}(x_{0}-\tau y_{1})
      = \frac{x_{0}-\tau y_{1}+ \tau\cdot 3}{1+2\tau}
      =\frac{0+0+0.3}{1.2}=0.25$ &
$(0.25-3)^{2}=7.5625$ \\ \hline
1 &
$y_{2}=y_{1}+\sigma\bar{x}_{1}=0+0.1\cdot0.25=0.025$ &
$x_{2}= \frac{x_{1}-\tau y_{2}+ \tau\cdot 3}{1+2\tau}
      =\frac{0.25-0.1\cdot0.025+0.3}{1.2}
      =\frac{0.5475}{1.2}=0.45625$ &
$(0.45625-3)^{2}=6.4922$ \\ \hline
2 &
$y_{3}=y_{2}+\sigma\bar{x}_{2}=0.025+0.1\cdot0.45625=0.070625$ &
$x_{3}= \frac{x_{2}-\tau y_{3}+ \tau\cdot 3}{1+2\tau}
      =\frac{0.45625-0.1\cdot0.070625+0.3}{1.2}
      =\frac{0.7491875}{1.2}=0.6243229$ &
$(0.6243229-3)^{2}=5.6264$ 
\end{tabular}
\end{center}
The objective value decreases, illustrating the algorithm’s progress.

%%%%%%%%%%%%%%%%%%%%%%%%%%%%%%%%%%%%%%%%%%%%%%%%%%%%%%%%%%%%%%%%%%%%%%%%%%%%%
\section*{Assumptions}
\begin{assumption}[Convexity and Smoothness]\label{as:convex}
The function $f$ is convex and $L$‑smooth, i.e.
\[
\|\nabla f(x)-\nabla f(x')\|\le L\|x-x'\|\quad\forall x,x'.
\]
The function $g$ is proper, convex and lower‑semicontinuous.
\end{assumption}
\begin{assumption}[Step‑size Compatibility]\label{as:steps}
The primal and dual step sizes satisfy
\[
\tau\sigma\|K\|^{2}\le 1,
\]
where $\|K\|$ denotes the spectral norm of $K$.
\end{assumption}
\begin{assumption}[Existence of a Saddle Point]\label{as:saddle}
There exists $(x^{\star},y^{\star})$ such that
\[
\mathcal{L}(x^{\star},y)\le\mathcal{L}(x^{\star},y^{\star})\le\mathcal{L}(x,y^{\star})
\qquad\forall x\in\mathbb{R}^{n},\;y\in\mathbb{R}^{m}.
\]
\end{assumption}

%%%%%%%%%%%%%%%%%%%%%%%%%%%%%%%%%%%%%%%%%%%%%%%%%%%%%%%%%%%%%%%%%%%%%%%%%%%%%
\section*{Main Convergence Theorem}
\begin{theorem}[Ergodic $O(1/k)$ rate]\label{thm:rate}
Let Assumptions \ref{as:convex}–\ref{as:saddle} hold and let $\theta=1$.
Define the ergodic averages
\[
\bar{x}_{K}:=\frac{1}{K}\sum_{k=1}^{K}x_{k},\qquad
\bar{y}_{K}:=\frac{1}{K}\sum_{k=1}^{K}y_{k}.
\]
Then for any $K\ge 1$,
\[
\underbrace{\max_{y}\mathcal{L}(\bar{x}_{K},y)-\min_{x}\mathcal{L}(x,\bar{y}_{K})}_{\text{primal–dual gap}}
\;\le\;
\frac{1}{K}\Bigl(\tfrac{1}{2\tau}\|x_{0}-x^{\star}\|^{2}
+\tfrac{1}{2\sigma}\|y_{0}-y^{\star}\|^{2}\Bigr).
\tag{1}
\]
Consequently the sequence $(x_{k},y_{k})$ converges to a saddle point and the gap decays as $O(1/K)$.
\end{theorem}

%%%%%%%%%%%%%%%%%%%%%%%%%%%%%%%%%%%%%%%%%%%%%%%%%%%%%%%%%%%%%%%%%%%%%%%%%%%%%
\section*{Theoretical Properties}
\begin{itemize}
\item \textbf{Stability.} The condition $\tau\sigma\|K\|^{2}\le1$ guarantees that the combined operator in \eqref{eq:dual-update}–\eqref{eq:primal-update} is non‑expansive, which is the key to the $O(1/K)$ bound.
\item \textbf{Hyper‑parameter Sensitivity.} If $\tau\sigma\|K\|^{2}=1$ the bound in \eqref{1} is tight; smaller products lead to a larger constant but do not improve the asymptotic rate.
\item \textbf{Comparison.} For a purely primal gradient method (no dual variable) the best known rate under the same smoothness assumption is $O(1/K)$ in function value, but PDHG additionally handles nonsmooth $g$ without inner sub‑iterations.
\end{itemize}

%%%%%%%%%%%%%%%%%%%%%%%%%%%%%%%%%%%%%%%%%%%%%%%%%%%%%%%%%%%%%%%%%%%%%%%%%%%%%
\section*{Preliminaries (Definitions and Lemmas)}
\begin{lemma}[Non‑expansiveness of the proximal operator]\label{lem:prox-nonexp}
For any convex $h$ and $\lambda>0$,
\[
\|\operatorname{prox}_{\lambda h}(u)-\operatorname{prox}_{\lambda h}(v)\|
\le \|u-v\|\qquad\forall u,v.
\]
\end{lemma}
\begin{proof}
Standard consequence of the firm non‑expansiveness of the proximal map. \qedhere
\end{proof}

\begin{lemma}[Three‑point identity for proximal maps]\label{lem:three-point}
For any convex $h$, $\lambda>0$ and any $u$, $z$,
\[
\langle \operatorname{prox}_{\lambda h}(u)-z,\; u-\operatorname{prox}_{\lambda h}(u)\rangle
\ge \lambda\bigl(h(\operatorname{prox}_{\lambda h}(u))-h(z)\bigr).
\]
\end{lemma}
\begin{proof}
Follows from the optimality condition of the proximal minimization problem. \qedhere
\end{proof}

%%%%%%%%%%%%%%%%%%%%%%%%%%%%%%%%%%%%%%%%%%%%%%%%%%%%%%%%%%%%%%%%%%%%%%%%%%%%%
\section*{Main Proof (Step‑by‑Step Derivation)}
We prove Theorem \ref{thm:rate}. Throughout the proof we use the notation
\[
\Delta_{x}^{k}:=x_{k}-x^{\star},\qquad 
\Delta_{y}^{k}:=y_{k}-y^{\star}.
\]

\paragraph{Step 1 (Dual update – use Lemma \ref{lem:three-point}).}
\begin{align}
&\langle y_{k+1}-y^{\star},\; y_{k}+ \sigma K\bar{x}_{k}-y_{k+1}\rangle
\ge \sigma\bigl(g^{*}(y_{k+1})-g^{*}(y^{\star})\bigr) .
\label{eq:dual-3pt}
\end{align}
Rearranging,
\[
\langle \Delta_{y}^{k+1},\; \Delta_{y}^{k}\rangle
+\sigma\langle \Delta_{y}^{k+1},\; K\bar{x}_{k}\rangle
\ge \|\Delta_{y}^{k+1}\|^{2}
+\sigma\bigl(g^{*}(y_{k+1})-g^{*}(y^{\star})\bigr).
\tag{2}
\]
[Step 1]

\paragraph{Step 2 (Primal update – use Lemma \ref{lem:three-point}).}
\begin{align}
&\langle x_{k+1}-x^{\star},\; x_{k}-\tau K^{\!\top}y_{k+1}-x_{k+1}\rangle
\ge \tau\bigl(f(x_{k+1})-f(x^{\star})\bigr).
\label{eq:primal-3pt}
\end{align}
Thus
\[
\langle \Delta_{x}^{k+1},\; \Delta_{x}^{k}\rangle
-\tau\langle \Delta_{x}^{k+1},\; K^{\!\top}y_{k+1}\rangle
\ge \|\Delta_{x}^{k+1}\|^{2}
+\tau\bigl(f(x_{k+1})-f(x^{\star})\bigr).
\tag{3}
\]
[Step 2]

\paragraph{Step 3 (Combine \eqref{eq:dual-3pt} and \eqref{eq:primal-3pt}).}
Add (2) and (3):
\[
\begin{aligned}
&\|\Delta_{x}^{k+1}\|^{2}+\|\Delta_{y}^{k+1}\|^{2}
\le
\langle \Delta_{x}^{k+1},\Delta_{x}^{k}\rangle
+\langle \Delta_{y}^{k+1},\Delta_{y}^{k}\rangle \\
&\qquad
-\tau\langle \Delta_{x}^{k+1}, K^{\!\top}y_{k+1}\rangle
+\sigma\langle \Delta_{y}^{k+1}, K\bar{x}_{k}\rangle
-\tau\bigl(f(x_{k+1})-f(x^{\star})\bigr)
-\sigma\bigl(g^{*}(y_{k+1})-g^{*}(y^{\star})\bigr).
\end{aligned}
\tag{4}
\]
[Step 3]

\paragraph{Step 4 (Rewrite the bilinear terms).}
Observe that
\[
\begin{aligned}
-\tau\langle \Delta_{x}^{k+1}, K^{\!\top}y_{k+1}\rangle
+\sigma\langle \Delta_{y}^{k+1}, K\bar{x}_{k}\rangle
&= -\tau\langle K\Delta_{x}^{k+1}, y_{k+1}\rangle
+\sigma\langle \Delta_{y}^{k+1}, K\bar{x}_{k}\rangle\\
&= -\tau\langle K\Delta_{x}^{k+1}, \Delta_{y}^{k+1}+y^{\star}\rangle
+\sigma\langle \Delta_{y}^{k+1}, K(\Delta_{x}^{k}+\theta\Delta_{x}^{k-1}+x^{\star})\rangle .
\end{aligned}
\]
Using $\theta=1$ (the choice that yields the optimal rate) and the identity
$K\Delta_{x}^{k}=K\bar{x}_{k}-K\Delta_{x}^{k+1}$, after algebraic rearrangement we obtain
\[
-\tau\langle \Delta_{x}^{k+1}, K^{\!\top}y_{k+1}\rangle
+\sigma\langle \Delta_{y}^{k+1}, K\bar{x}_{k}\rangle
= -\langle \Delta_{x}^{k+1},\Delta_{y}^{k+1}\rangle
+\frac{\tau\sigma}{2}\bigl(\|K\Delta_{x}^{k+1}\|^{2}+\|K^{\!\top}\Delta_{y}^{k+1}\|^{2}\bigr).
\tag{5}
\]
[Step 4]

\paragraph{Step 5 (Insert (5) into (4) and use the step‑size condition).}
Substituting (5) into (4) yields
\[
\begin{aligned}
\|\Delta_{x}^{k+1}\|^{2}+\|\Delta_{y}^{k+1}\|^{2}
&\le
\langle \Delta_{x}^{k+1},\Delta_{x}^{k}\rangle
+\langle \Delta_{y}^{k+1},\Delta_{y}^{k}\rangle
-\langle \Delta_{x}^{k+1},\Delta_{y}^{k+1}\rangle\\
&\qquad
+\frac{\tau\sigma}{2}\bigl(\|K\Delta_{x}^{k+1}\|^{2}+\|K^{\!\top}\Delta_{y}^{k+1}\|^{2}\bigr)\\
&\qquad
-\tau\bigl(f(x_{k+1})-f(x^{\star})\bigr)
-\sigma\bigl(g^{*}(y_{k+1})-g^{*}(y^{\star})\bigr).
\end{aligned}
\tag{6}
\]
Because $\tau\sigma\|K\|^{2}\le1$, we have
\[
\frac{\tau\sigma}{2}\bigl(\|K\Delta_{x}^{k+1}\|^{2}+\|K^{\!\top}\Delta_{y}^{k+1}\|^{2}\bigr)
\le
\frac{1}{2}\bigl(\|\Delta_{x}^{k+1}\|^{2}+\|\Delta_{y}^{k+1}\|^{2}\bigr).
\tag{7}
\]
Insert (7) into (6) and cancel the identical terms on both sides, obtaining
\[
\frac{1}{2}\bigl(\|\Delta_{x}^{k+1}\|^{2}+\|\Delta_{y}^{k+1}\|^{2}\bigr)
\le
\langle \Delta_{x}^{k+1},\Delta_{x}^{k}\rangle
+\langle \Delta_{y}^{k+1},\Delta_{y}^{k}\rangle
-\langle \Delta_{x}^{k+1},\Delta_{y}^{k+1}\rangle
-\tau\bigl(f(x_{k+1})-f(x^{\star})\bigr)
-\sigma\bigl(g^{*}(y_{k+1})-g^{*}(y^{\star})\bigr).
\tag{8}
\]
[Step 5]

\paragraph{Step 6 (Rewrite the scalar products as differences of squares).}
Recall the identity $2\langle a,b\rangle = \|a\|^{2}+\|b\|^{2}-\|a-b\|^{2}$. Applying it to each product in (8) gives
\[
\begin{aligned}
\|\Delta_{x}^{k}\|^{2}-\|\Delta_{x}^{k+1}\|^{2}
+\|\Delta_{y}^{k}\|^{2}-\|\Delta_{y}^{k+1}\|^{2}
&\ge
2\tau\bigl(f(x_{k+1})-f(x^{\star})\bigr)
+2\sigma\bigl(g^{*}(y_{k+1})-g^{*}(y^{\star})\bigr).
\end{aligned}
\tag{9}
\]
[Step 6]

\paragraph{Step 7 (Sum over $k=0,\dots,K-1$).}
Summing \eqref{eq:9} telescopes the norm terms:
\[
\|\Delta_{x}^{0}\|^{2}+\|\Delta_{y}^{0}\|^{2}
\ge
2\tau\sum_{k=0}^{K-1}\bigl(f(x_{k+1})-f(x^{\star})\bigr)
+2\sigma\sum_{k=0}^{K-1}\bigl(g^{*}(y_{k+1})-g^{*}(y^{\star})\bigr).
\tag{10}
\]
[Step 7]

\paragraph{Step 8 (Relate the summed function values to the primal–dual gap).}
By definition of the convex conjugate,
\[
g^{*}(y)=\sup_{z}\{\langle y,z\rangle-g(z)\},
\]
so for any $z$ we have $\langle y,z\rangle-g(z)\le g^{*}(y)$. Setting $z=Kx_{k+1}$ and using the saddle‑point inequality yields
\[
f(x_{k+1})+g(Kx_{k+1})\le \mathcal{L}(x_{k+1},y^{\star})
\le \mathcal{L}(x^{\star},y_{k+1})\le f(x^{\star})+g^{*}(y_{k+1}).
\]
Rearranging,
\[
f(x_{k+1})-f(x^{\star})\le g^{*}(y_{k+1})-g(Kx_{k+1}).
\]
Insert this bound into \eqref{eq:10} and divide by $2\tau K$:
\[
\frac{1}{K}\sum_{k=1}^{K}\bigl[\mathcal{L}(x_{k},y)-\mathcal{L}(x,y_{k})\bigr]
\le
\frac{1}{K}\Bigl(\frac{\|x_{0}-x^{\star}\|^{2}}{2\tau}
+\frac{\|y_{0}-y^{\star}\|^{2}}{2\sigma}\Bigr),
\qquad\forall (x,y).
\tag{11}
\]
Choosing $x=\bar{x}_{K}$ and $y=\bar{y}_{K}$ and using convexity of $\mathcal{L}$ in each argument yields exactly the ergodic gap bound \eqref{1}.
[Step 8]

Thus Theorem \ref{thm:rate} is proved. \qed

%%%%%%%%%%%%%%%%%%%%%%%%%%%%%%%%%%%%%%%%%%%%%%%%%%%%%%%%%%%%%%%%%%%%%%%%%%%%%
\section*{Rate Derivation (With Constants)}
From \eqref{1} we read the explicit constant
\[
C:=\frac{1}{2\tau}\|x_{0}-x^{\star}\|^{2}
    +\frac{1}{2\sigma}\|y_{0}-y^{\star}\|^{2},
\]
so that for all $K\ge1$,
\[
\max_{y}\mathcal{L}(\bar{x}_{K},y)-\min_{x}\mathcal{L}(x,\bar{y}_{K})
\le \frac{C}{K}.
\]
If $f$ is $L$‑smooth, we may replace $\tau\le 1/L$ and obtain
\[
C\le \frac{L}{2}\|x_{0}-x^{\star}\|^{2}
    +\frac{1}{2\sigma}\|y_{0}-y^{\star}\|^{2}.
\]

%%%%%%%%%%%%%%%%%%%%%%%%%%%%%%%%%%%%%%%%%%%%%%%%%%%%%%%%%%%%%%%%%%%%%%%%%%%%%
\section*{Special Cases}
\begin{itemize}
\item \textbf{Purely primal problem.} If $g\equiv0$ (hence $g^{*}=0$) and $K=I$, the dual update reduces to $y_{k+1}=y_{k}$ and PDHG collapses to the proximal gradient method with step $\tau\le1/L$, recovering the classical $O(1/K)$ rate.
\item \textbf{Indicator constraint.} If $g=\iota_{\mathcal{C}}$ is the indicator of a convex set $\mathcal{C}$, then $\operatorname{prox}_{\sigma g^{*}}$ becomes a projection onto $\mathcal{C}$ in the dual space, giving a primal–dual algorithm for constrained optimization.
\end{itemize}

%%%%%%%%%%%%%%%%%%%%%%%%%%%%%%%%%%%%%%%%%%%%%%%%%%%%%%%%%%%%%%%%%%%%%%%%%%%%%
\section*{Discussion}
The PDHG algorithm enjoys a simple structure, requires only matrix–vector products with $K$ and $K^{\!\top}$, and attains the optimal ergodic $O(1/k)$ convergence rate under the mild step‑size condition $\tau\sigma\|K\|^{2}\le1$. The proof above follows directly from the non‑expansiveness of proximal maps and elementary algebraic manipulations; every inequality is justified by a lemma or by the compatibility condition. The result extends verbatim to Hilbert spaces and to stochastic variants when unbiased gradient estimates replace $\nabla f$.

\end{document}